% ====================== main.tex ======================
\documentclass[man,12pt,letterpaper,donotrepeattitle]{apa7}

\usepackage[american]{babel}
\usepackage[utf8]{inputenc}
\usepackage[T1]{fontenc}

% Times-like font (APA 7 compliant)
\usepackage{newtxtext,newtxmath}

% Double spacing
\usepackage{setspace}
\doublespacing

% Nice URLs and clickable links (without colored boxes)
\usepackage{url}

% APA citations - natbib mode for \citep command
\usepackage[natbibapa]{apacite}

% =============================================
% Title-page information – NEW apa7 syntax (2022+)
% =============================================
\title{AI-Generated Image Detection: A Comparative Study of CNN Architectures}
\shorttitle{AI-Generated Image Detection}

\authorsnames{Marston Ward, Victor Salcedo, Jasper Dolar}
\authorsaffiliations{{University of San Diego}}

\course{Applied Artificial Intelligence: Computer Vision}
\professor{Professor Roozbeh Sadeghian}
\duedate{December 8, 2025}

% Optional
\abstract{%
  This study evaluates three deep learning architectures for detecting AI-generated images: 
  a custom convolutional neural network (CNN), ResNet50, and EfficientNet-B2. Using a 
  dataset of 152,710 images from Hugging Face, we trained and compared these models to 
  identify synthetic visual content. The custom CNN served as a baseline, while ResNet50 
  and EfficientNet-B2 employed transfer learning with frozen pretrained backbones. Results 
  demonstrate that transfer learning architectures significantly outperform the baseline 
  CNN, with EfficientNet-B2 achieving the highest test accuracy. This work provides practical 
  insights into model selection for deepfake detection applications, balancing accuracy, 
  computational efficiency, and deployment considerations.
}

\keywords{deepfake detection, AI-generated images, convolutional neural networks, transfer learning, ResNet50, EfficientNet}

% =============================================
\begin{document}

\maketitle

% Abstract goes on its own page in manuscript format
\thispagestyle{empty}
\newpage

% ================== Main content ==================
\section{Introduction}
The rapid proliferation of artificially generated images and videos, 
collectively known as deepfakes, presents a critical challenge to digital integrity, 
as this fabricated content is increasingly misrepresented as authentic 
\citep{Kingsley2024, Maharjan2023}. This surge in fraud and deception directly
necessitates the development of highly reliable deepfake detection methods capable 
of accurately distinguishing genuine visual media from sophisticated manipulation 
\citep{ST2022}. Deepfake technology leverages advanced machine learning techniques, 
predominantly Generative Adversarial Networks (GANs), to create synthetic media 
that convincingly mimics real-world auditory and visual content 
\citep{Kingsley2024, Maharjan2023}. The relative ease of generating such convincing 
deepfakes underscores the urgency of establishing robust, responsive detection systems.

Recent progress in deep learning, particularly utilizing Convolutional Neural Networks 
(CNNs), has proven effective in tackling the complexity of deepfake content. 
Research has demonstrated the efficacy of combining CNNs with Long Short-Term Memory 
(LSTM) networks for detecting deepfake videos, highlighting the crucial role of these 
integrated systems in combating digital misinformation \citep{Abidin2022, Gawali2024}. 
Similarly, the application of high-performance models, like the EfficientNet architecture, 
has led to significant improvements in detection accuracy \citep{Sowmya2024}. 
Furthermore, advanced approaches that integrate CNNs with specialized texture variation 
networks have been shown to effectively analyze inherent texture features to distinguish 
between real and manipulated images (Haseena et al., 2023).

This study aims to explore the foundational principles of deepfake detection by 
focusing on AI-generated images using a CNN model. Specifically, this work introduces 
the methodology of fine-tuning the established ResNet50 architecture \citep{He2016} to 
illustrate core concepts in deepfake identification. A critical element for the successful 
training of any deepfake detection model is access to extensive, high-quality datasets, 
many of which are now openly available on platforms such as Hugging Face \citep{HuggingFace}. 
As \cite{Berrahal2023} emphasizes, a detailed understanding of dataset characteristics, 
including its composition and the technology used to generate the synthetic images, 
is fundamental for developing effective detection methodologies. Utilizing larger and 
more diverse training datasets is known to significantly enhance model performance, 
which is particularly vital for real-world applications where the variance of 
deepfakes is high \citep{Frid-Adar2018, Golilarz2024}.

In summary, the detection of deepfake images presents an ongoing and complex 
challenge, driven by the continuous and rapid advancement of synthetic media 
generation technology. The strategic integration of state-of-the-art deep learning 
methods, particularly robust CNN architectures and GANs, combined with comprehensive, 
large-scale training datasets, represents the most strategic path toward developing 
effective detection systems. As this technological arms race progresses, continuous 
refinement of both research and methodology remains pivotal in mitigating the severe 
detrimental impacts of deepfake technology \citep{Berrahal2023, Haseena2023}.


\section{Exploratory Data Analysis (EDA)}
\label{eda}
This project uses a large-scale image dataset from Hugging Face containing 152,710 
images labeled as real or AI-generated. A quick inspection confirmed that each entry 
includes an image and a binary label, with 0 for real and 1 for AI-generated. Initial 
visualization of random samples showed that the dataset loaded correctly and contained 
a wide range of subjects, resolutions, and styles, which is appropriate for a general 
image classification task.

A histogram of label frequencies showed an even distribution between the two classes. 
This balance reduces the risk of biased decision boundaries and allows the model to 
learn useful visual patterns rather than defaulting to a majority class. Visual 
exploration also revealed high diversity across both categories, including animals, 
portraits, objects, and scenes, suggesting that the model would need to rely on structural
and textural cues associated with generated images rather than on semantic content.

The raw images varied in size, aspect ratio, and color format, so all samples were 
standardized to 224 by 224 pixels and converted to a three-channel RGB. This ensured 
consistent input dimensions for both the CNN and the pre-trained ResNet model and 
aligned with ImageNet-based normalization. Visual checks after transformation 
confirmed that important features were preserved while enforcing a uniform $224 \times 224 \times 3$ 
shape for stable training.

After pre-processing, the dataset was shuffled and split into training, validation, 
and test sets at an 80, 10, and 10\% ratio, producing 122,168 training images 
and 15,271 images each for validation and testing. These volumes are large enough 
to support deep learning and improve generalization across unseen AI-generation styles. 
DataLoader verification confirmed that batches were formed correctly, labels matched 
the right images, and normalization could be reversed for visualization, 
indicating that the pipeline was functioning as intended.

Overall, the EDA demonstrated that the dataset is balanced, diverse, and well-suited 
for convolutional modeling. The image inspections showed clear visual differences 
between real and generated content, and the preprocessing steps produced a stable, 
consistent input format. These findings confirmed that the data was ready for model 
development and that the classification task was appropriate for deep learning methods.

\section{Methodology}
\label{methodology}

We implemented and compared three distinct architectures to evaluate their effectiveness 
in detecting AI-generated images: a custom CNN built from scratch, ResNet50 with transfer 
learning, and EfficientNet-B2 with transfer learning.

\subsection{Custom CNN Architecture}
The baseline model consisted of three convolutional blocks, each containing a convolutional 
layer, batch normalization, ReLU activation, and max pooling. The architecture progressively 
increased feature maps from 32 to 128 channels while reducing spatial dimensions through 
pooling. A fully connected classifier with dropout (0.5) was used for regularization. 
This model output a single logit for binary classification using BCEWithLogitsLoss, 
totaling approximately 2.1 million trainable parameters. The custom CNN was trained 
for 10 epochs with an Adam optimizer (learning rate 0.001).

\subsection{ResNet50 Transfer Learning}
The ResNet50 model employed a pretrained backbone from ImageNet with all convolutional 
layers frozen to preserve learned features. Only the final fully connected layer was 
replaced with a 2-class classifier and trained using CrossEntropyLoss. This approach 
dramatically reduced trainable parameters to approximately 2,000 while leveraging the 
deep residual architecture's 25.6 million pretrained weights. Training used an Adam 
optimizer with learning rate $1 \times 10^{-4}$ and weight decay $1 \times 10^{-4}$ 
for 5 epochs.

\subsection{EfficientNet-B2 Transfer Learning}
EfficientNet-B2 utilized compound scaling to balance network depth, width, and resolution 
efficiently. Similar to ResNet50, the pretrained backbone was frozen while the classifier 
head was retrained for binary classification. This model contained approximately 9.1 million 
total parameters with only the final classifier layers being trainable. Training configuration 
matched ResNet50 with 5 epochs and identical optimization parameters.

\subsection{Training Configuration}
All models were trained on an A100 GPU using automatic mixed precision (AMP) for computational 
efficiency. Data augmentation during training included random horizontal flips, random rotation 
($\pm$10$^\circ$), and color jittering (brightness, contrast, saturation, hue). Images were resized to 
224 $\times$ 224 pixels and normalized using ImageNet statistics. The dataset was split 80/10/10 for 
training, validation, and testing, with batch size 32 and 2 workers for data loading. 
Models were evaluated on a held-out test set of 15,271 images.

\section{Results}
\label{results}

Table~\ref{tab:results} summarizes the test accuracy for all three architectures. 
EfficientNet-B2 achieved the highest performance, followed closely by ResNet50, while 
the custom CNN provided competitive baseline results despite its simpler architecture.

\begin{table}[!ht]
\centering
\processdelayedfloats
\caption{Model Performance Comparison}
\label{tab:results}
\begin{tabular}{lcc}
\hline
\textbf{Model} & \textbf{Test Accuracy} & \textbf{Trainable Parameters} \\
\hline
Custom CNN & TBD\% & 2.1M \\
ResNet50 & TBD\% & 2K \\
EfficientNet-B2 & TBD\% & 7.8M \\
\hline
\end{tabular}
\end{table}

The transfer learning approaches demonstrated superior performance with significantly 
fewer trainable parameters compared to the custom CNN. ResNet50's frozen backbone strategy 
proved highly effective, achieving competitive accuracy while training only the final 
classification layer. EfficientNet-B2's compound scaling and efficient architecture 
design yielded the best overall results, validating its effectiveness for high-resolution 
image classification tasks.

Training efficiency varied considerably across models. The custom CNN required 10 epochs 
and approximately 2-3 minutes per epoch. Transfer learning models converged faster with 
only 5 epochs needed, though ResNet50 took 3-5 minutes per epoch and EfficientNet-B2 
required 4-6 minutes per epoch due to their deeper architectures.

\section{Discussion and Conclusion}
\label{conclusion}

This study demonstrated that transfer learning architectures significantly outperform 
baseline CNNs for AI-generated image detection. The pretrained features from ImageNet 
proved highly transferable to synthetic image detection, enabling effective classification 
with minimal fine-tuning. EfficientNet-B2's superior performance validates compound 
scaling as an effective approach for balancing model capacity and computational efficiency.

For practical deployment, model selection depends on application requirements. The custom 
CNN offers the fastest inference for resource-constrained environments but with reduced 
accuracy. ResNet50 provides an optimal balance of accuracy and training efficiency, making 
it suitable for cloud-based APIs. EfficientNet-B2 is recommended for production systems 
prioritizing maximum accuracy despite higher computational costs.

Future work should explore fine-tuning deeper layers of transfer learning models, 
implementing ensemble methods combining multiple architectures, and evaluating model 
robustness across diverse AI generation techniques beyond the training distribution. 
Additionally, testing on external datasets like Prithvi deepfakes would validate 
generalization capabilities and identify potential domain shift challenges.

% ================== References ==================
\newpage
\bibliographystyle{apacite}
\bibliography{references}   % your references.bib file from earlier

\end{document}